\section{Kublai Khan}

\begin{quotex}
This is a chapter from \textit{The Middle Ages} by \textbf{Karl Heyer}.
\end{quotex}

Just as there was only one Ever Blue Sky and one Great Spirit, so there was to be only one ruler on earth, chosen by the Great Spirit. It was an attitude expressed in Genghis Khan's impressive formulation: One divine ruler in heaven and one Great Khan to rule the earth!

The was the spiritual foundation for the Mongolian desire for uniformity and power. It its service they placed the might of their Mars qualities, their organisational talent and their capacity to maintain a great state. All this was connected with the lack among nomads peoples of any cultural creativity, a lack that had already been visible in the great Turkic empire of the fifth century.

Exemplary order reigned within Genghis Khan's empire: an end was put to robbery along the great trade routes, there was a postal system using mounted carriers, and trade developed freely. The ``Pax Mongolica'' covered an immense area. At the death of Genghis Khan and beyond, his realm stretched from Korea and the Yellow Sea to the Black and Arabian Seas. Communications were highly developed with a road network covering the whole empire that enabled the cultures of the east and west to meet.

The overcoming of distance achieved by European technology at the beginning of the twentieth century was anticipated in the thirteenth by the wild will forces of Genghis Khan and taken to its ultimate perfection in Kublai Khan's imperial organisation. Merchants using the route from the Genoese colony of Tanais on the shores of the Sea of Azov to China reported that it was entirely safe by day and night as a result of the ``Pax Mongolica'' or ``Pax Tatarica''.

Thus the thirteenth century saw world commerce made possible by the unnatural, superhuman forces of expansion developed by a primitive people on horseback. It was not to last long, however, and by the second half of the fourteenth century the unity of the European and Asian continent had collapsed once more.

As mentioned above, the organisation of the Mongolian state reached its culmination in the realm of Kublai Khan (d. 1294 ) who was a grandson and successor of Genghis Khan. He ruled China and the Far East while the middle and western regions had already become relatively independent under other descendants of Genghis Khan.

Rudolf Steiner stated that Kublai Khan, too, was still influenced by the initiation discussed in connection with Genghis Khan. At his brilliant court resided for many years the Venetian \textbf{Marco Polo} (1254-1324) whose book about ``The Wonders of the World'' for the first time described this remote Chinese-Mongolian realm. For a long time, his tales met with disbelief, so fantastic and astonishing did they sound.

Kublai Khan had moved his residence from Mongolia to Peking, and his realm now achieved something like a synthesis of Mongolian organisational energy and talent with the culture of the Chinese. This entailed a watering down of the wilder aspects of the Mongolian character and this ultimately a part in bringing about the end of the Mongolian rule in China followed by the dissolution of the Mongolian empire as a whole.

Meanwhile, however, conditions in Kublai Khan's realm were truly remarkable in that they showed what a synthesis between the wild nature of nomadic Mongolians with the high culture of those of China was capable of achieving.

Descriptions of Kublai's realm reveal an exemplary administration, far-reaching social welfare and many other characteristics of a flourishing community full of social concern. Joachim Barckhausen considered that, measured against its power and extent, its technical and social institutions need not appear so very astonishing. From the vantage point of the twentieth century, it not ought to be too difficult to gain an idea of the capacities of a totalitarian state that makes full use of all its powers.

The court at Peking receives enormous revenues and functions as the heart of the realm, quickening its pulse by the exchange of goods. Anyone becoming poor through no fault of his own receives state support equal to his normal standard of living. Soup kitchens cater for the poorest classes. There is an obligation to provide labour for the state.

Wherever there is a question of organisation, the Mongolians are exemplary. They bring a kind of Asian Prussianess to the Far East. With the help of an efficient policing system criminality is appreciably reduced.

This ideal, realised by the Mongolian Chinese state under Kublai Khan, reveals nuances that remind us of the modern Chinese writer Ku Hung Ming who described the significant elements of Chinese administration in connection with its ancient ideals. Ku Hung Ming compared the western Christian question: ``What is the highest purpose of the human being'' with that of the Confucian catechism: ``What is the highest purpose of the citizen?''

For him the highest purpose of the individual as a moral being is ``to be an obedient son and a good citizen''. He compared the Chinese person with a domestic animal, tame and lacking in individuality and thus easily led. A Chinese is an eminently social creature because the impulse to individuality has not yet dawned in him.

In all seriousness Ku Hung Ming therefore suggested that in order to escape from the strife of the western world people ought to be changed into true Chinese. Repeatedly and emphatically he challenged the western nations to tear up their Magna Carta of freedom and establish in its place ``a Magna Carta of loyalty such as we Chinese possess in our religion of the good citizen''.

All the impressive social structures of eastern Asia have grown and developed on this basis of old pre -- individual sense of community. The individual is embedded in the community as though in a large family. He is easily led (as at the more Chinese, passive end of the scale) and ready to make sacrifices on behalf of the community (as at the more active end of the scale seen in the Japanese or the Tartar Mongolians). These community structures take absolutely no account of the individuality becoming aware of itself, and the western problem of the tension between the individual and the community is non -- existent as far as they are concerned.

We have already mentioned that the Mongolians of Genghis Khan and his successors lost their vitality and perished as a result of their intense contact with their cultured neighbours. Finally the Chinese rose up successfully against their Mongolian overlords.

With a proclamation beginning with the words: ``These barbarians are created to obey and not to command a cultured nation''. China began the attack, therefore taking only a few months to shake off the yoke. In 1388 the Mongolians were beaten in Mongolia by a Chinese army. They relinquished the proud name of Mongolians given to them by Genghis Khan and henceforth once more called themselves Tatars.

Bloody feuds had already been going on for some time among the other more westerly tribes of the Mongolian empire, especially between the Persian Khans and the Russian Golden Horde. These feuds absorbed much energy and it is probably to them in the first place that Europe owes the fact that the Mongolians did not return after the Battle of Legnica in 1241 or come back at any later time. Their domination never stretched further than their Russian realm.

The dissolution of the Mongolian realm took place in silence, almost uncannily. Hollowed out as though by termites the colossus still stood for a number of decades merely because in an Asia weakened by what had passed there was nothing to take its place. Finally, almost without any external inducement, it collapsed and fell to dust.

Something that points to deep secrets of human evolution is the fact that Buddhism found its way to the peoples we call Mongolian in the widest sense. This occurred in China in the first century A.D. in Tibet in the seventh, and in Japan in the sixth. The Mongolian Tatars, however, only adopted Buddhism after the time of Genghis Khan.

Rudolf Steiner explained the spiritual realities lying behind this fact when he spoke of Buddha as having been the bearer of a divine entity who had been incarnated among human beings during the time of Atlantis when human bodies were not yet as dense as they later became. Because this was so, human beings came to know gods such as Thor, Odin, Baldur and Zeus as companions on earth. Later, when the human form became denser as it descended more deeply into the material world, these beings remained in the supersensible world and became invisible to human eyes.

\begin{quotationx}
Those human beings who purified their more delicate bodies, by living in accordance with the instructions given by the initiates, went out in a certain way towards the gods. Thus, although they were incarnated in fleshly bodies, if they purified themselves, they became capable of being overshadowed by a spiritual being who could not descend as far as a physical body. In this way Buddha became a vessel for Odin. The god named Odin in Germanic mythology appeared again as Buddha.

Thus it is possible to say that much of what was secret in Atlantean times passed over into the message proclaimed by Buddha. What Buddha experienced was in harmony with what the gods had experienced in those spiritual realms and what human beings themselves had experienced when they still lived in those realms. When Odin's teachings reappeared they took little account of the physical plane apart from stressing that it is an abode of pain from which it is important to seek release.

Much of what had been in Odin's being spoke through Buddha, and that is why Buddha's teachings met with the deepest understanding in those who were latecomers from Atlantis. Among the Asiatic population there are peoples who have definitely remained behind at the Atlantean stage, although outwardly, of course, they have kept pace with earthly evolution. In the Mongolian peoples much remains that comes from Atlantis; they are stragglers following on from the ancient population of Atlantis. The more stationery stream of the Mongolian population represents an inheritance from Atlantis and that is why the teachings of Buddha serve these peoples particularly well and why Buddhism has gained so much ground among them. 
\flright{\emph{Egyptian Myths and Mysteries}}
\end{quotationx}

The significance of this fact is that the bellicose, aggressive, Mongolian race of Mars, as though led by wise spiritual forces, has found a link with the Buddha, a being who is the very opposite of aggressive and warlike and instead emanates calm peacefulness, inwardness, compassion and love towards human beings. Far from being Mars-like, the Buddha is entirely Mercury-like.

It is as though his influence were intended to provide a healing antithesis for the Mongolian peoples. Note: Sigismund von Gleich points out the surprising contrast between the peaceful, almost feminine traits in the Chinese soul and the characteristics of the Turanian-Mongolian races. Perhaps this contrast indicates that very early on a ``Mercurial'' counter-impulse went to work on the Martian, Turanian characteristics in the Chinese. This would also help us to understand another characteristic of the Chinese to which von Gleich also points, namely, their affinity with calculation and mathematical thinking on which their talent for commercial matters is founded. This places the phenomenon of Mongolian acceptance of Buddha in the context of the great link that exists in world evolution between Mars and Mercury. Mars must be transformed into Mercury; the effects of Mars must be ``healed'' by Mercury. This belongs to the mission of Buddha as it was later fulfilled as a deed of redemption in the supersensible Mars sphere.

The Mongolians of Genghis Khan adhered to the religion and magic of Shamanism in which decadent Atlantean forces were at work and which threatened the western world as a result of the Mongolian invasions. It is understandable that Genghis Khan himself remained faithful to these beliefs until the end of his life.

The final years of this wild conqueror, however, appeared to have been surrounded in a strange way by the legends of Buddhism and Tibetan Lamaism. He approached the regions where they held sway in his campaign against the kingdom of Hsia, the northeastern threshold to the highlands of Tibet. According to a Mongolian chronicle, Genghis Khan captured the daughter of the king of Hsia, a kind of Buddhist magician, priest and robber:

\begin{quotex}
He took her with him into his yurt, but she had learned the secret arts of the Tibetans from her father. As he lay with her for the first night, she did a magic harm to Genghis Khan that henceforth ailed him.
\end{quotex}

In retelling this story recorded by a Mongolian prince of a later century who had meanwhile converted to Buddhism, Barckhausen pointed out its obvious symbolism: Tibet took revenge for the sacrilege committed against it and slew Genghis Khan. ``Some day in the future, when the yellow priests descend from the forbidden mountains, the whole of the Mongolian nation will suffer what their Khan suffered.''

Expressing the great historical phenomenon in his own way, Barckhausen continued: ``The religion of Buddha has turned the Mongolians into powerless, defenseless devils. Because of it the once mightiest and strongest nation on earth is threatened with extinction today.''


\flrightit{Posted on 2022-08-29 by Mu Nianci}

